\documentclass[../main.tex]{subfiles}

\begin{document}
    Berdasarkan hasil praktikum pada Modul {\thechapter} tentang {\leftmark}, dapat diambil kesimpulan bahwa:
    \begin{enumerate}
        \item {\Pointer} adalah variabel yang menyimpan alamat memori dari variabel lain. Pointer dalam C++ dapat dibuat dengan cara menaruh tanda \citecode{*} di depan tipe data pada saat inisialisasi variabel.
        \item Ada perbedaan penting antara \textit{struct of pointer} dan \textit{pointer of struct}. \textit{Struct of pointer} mengacu pada struktur data yang berisi {\pointer}, sementara \textit{pointer of struct} adalah {\pointer} yang menunjuk ke suatu struktur data.
        \item Penggunaan {\pointer} memungkinkan pengalokasian dan manipulasi memori secara dinamis. Dengan {\pointer}, memori dapat dialokasikan, digunakan, dan dibebaskan sesuai kebutuhan program, menggunakan fungsi-fungsi seperti \citecode{malloc}, \citecode{calloc}, \citecode{realloc} dan \citecode{free} dalam bahasa C. Dua kata kunci \citecode{new} dan \citecode{delete} ditambahkan pada bahasa C++ untuk mengalokasikan dan membebaskan memori.
        \item Data dapat diterima oleh fungsi dengan cara menggunakan parameter \textit{pass-by} {\pointer} dan \textit{pass-by} {\reference}. Variabel yang dimasukan sebagai argumen fungsi yang menerima alamat memori variabel bukan nilai variabel.
    \end{enumerate}
\end{document}
